\documentclass[11pt]{article}

% =====================================================================
% Internal APIs Are All You Need  (arXiv:2604.00694v1)
% Shadow APIs, Shared Discovery, and the Case Against Browser-First
% Agent Architectures. Paper I of the trilogy — the wedge. Academic
% format mirrors the arXiv source.
% =====================================================================

\usepackage[margin=1in]{geometry}
\usepackage{hyperref}
\usepackage{amsmath}
\usepackage{amssymb}
\usepackage{enumitem}
\usepackage{booktabs}
\usepackage{xcolor}
\hypersetup{colorlinks=true, linkcolor=blue, urlcolor=blue, citecolor=blue}

\title{\textbf{Internal APIs Are All You Need}\\
\large Shadow APIs, Shared Discovery, and the Case Against\\
Browser-First Agent Architectures}

\author{Lewis Tham\thanks{Unbrowse AI, \texttt{lewis@unbrowse.ai}.}
  \and Nicholas Mac Gregor Garcia\thanks{School of Computing, National University
  of Singapore, \texttt{ngarcia@nus.edu.sg}.}
  \and Jungpil Hahn\thanks{School of Computing, National University of Singapore,
  \texttt{jungpil@nus.edu.sg}.}}
\date{April 2026 \quad(arXiv:2604.00694v1)}

\begin{document}
\maketitle

\begin{abstract}
Autonomous agents increasingly interact with the web, yet most websites remain
designed for human browsers: an agent that wants a single fact or action loads a
full page, executes its scripts, and waits on late requests to recover a call the
site's own front-end already knew how to make. We observe that this cost is paid
over and over for information that does not change between callers, and that the
structured request behind a page --- a \emph{shadow API}, the first-party endpoint
the front-end calls --- is learnable once and replayable by anyone. We present
Unbrowse, a system that passively learns callable routes from real usage and stores
them in a shared index, so that an intent resolves to a ranked endpoint shortlist
and executes against the live site without a browser. We frame the economics as a
single adoption condition --- a cost-minimising agent reuses a shared route if and
only if the route fee undercuts the cost of rediscovering it --- and settle paid
execution over x402 so the indexer who mapped a route, the platform, and the site
owner can be compensated in one transaction. Across 94 live domains we measure a
\textbf{3.6$\times$ mean} and \textbf{5.4$\times$ median} speedup over a
browser-automation baseline, with a best single-domain speedup of $30\times$. The
short version: when a website already exposes a stable first-party route behind
its UI, that route should be the agent's first interface and the browser should be
the fallback.
\end{abstract}

\section{Introduction}
The web was built for humans to navigate pages. Agents do not want pages --- they
want actions. When an autonomous agent is asked for the price of a product, the
status of an order, or the results of a search, it does not need the page's markup,
its analytics scripts, or its cookie banner; it needs the one structured request
the page makes to its own backend, and the structured response that comes back.
Browser-first agent architectures recover that request the expensive way: they
render the whole page, run its JavaScript, and observe the resulting network
traffic, paying a full browser's latency and token cost on every call.

We make one claim and try not to oversell it: for many repeatable agent web
tasks, the browser is unnecessary overhead once the first-party endpoint behind
the page --- the \emph{shadow API} --- has been learned and verified fresh. The
discovery cost should be paid once per route and amortised across later callers,
not re-paid on every visit. This paper presents
Unbrowse, a system built on that observation, and makes three contributions:
\begin{enumerate}[leftmargin=1.6em]
  \item A passive route-discovery pipeline that learns callable interfaces from
  real usage and normalises them into a shared graph.
  \item An adoption condition that states precisely when a shared route commons is
  the cost-minimising choice, and a three-tier payment architecture that settles
  the resulting value over x402.
  \item An evaluation across 94 live domains showing a 3.6$\times$ mean and
  5.4$\times$ median speedup over a browser baseline, with the reuse mechanism
  isolated in a runnable reference benchmark.
\end{enumerate}

\section{Related work}
\subsection{Web agents and browser automation}
A large body of work has agents drive real browsers and act on rendered pages:
end-to-end multimodal web agents~\cite{webvoyager}, and benchmarks such as
Mind2Web~\cite{mind2web} and WebArena~\cite{webarena} that measure agents on live
or simulated sites. This line treats the rendered page as the interface. We take
the opposite stance: the rendered page is a fallback, and the durable interface is
the first-party API the page already calls.

\subsection{Tool learning and API discovery}
Toolformer teaches a model \emph{when} to call an API~\cite{toolformer} and Gorilla
teaches it to emit \emph{correct} calls across many APIs~\cite{gorilla}. These
solve single-agent tool use against \emph{declared} interfaces. They do not address
\emph{undeclared} first-party endpoints, nor what happens when many agents
rediscover the same interface repeatedly --- the shared-discovery problem we take
up here.

\subsection{API aggregation platforms}
Aggregators and standard connectors expose curated, declared APIs behind a uniform
surface. Their coverage is bounded by what providers choose to publish. Shadow-API
discovery is complementary and unbounded in a different direction: it learns the
endpoints a site already serves to its own front-end, whether or not they are
documented.

\subsection{Agent economies and communication protocols}
The Model Context Protocol standardises how a model connects to tools and
data~\cite{mcp}; x402 revives HTTP 402 as a real micropayment
rail~\cite{x402}. Prior agent-payment work largely stops at ``the agent can pay.''
We extend it to ``the payment routes to everyone who created the value'' and frame
the commons in the tradition of Ostrom's work on common-pool
resources~\cite{ostrom}.

\section{System design}
\subsection{System architecture}
Unbrowse resolves an intent along a three-path model: a cache-first lookup against
the shared route graph, a graph traversal over typed endpoint edges, and --- only
on a miss --- a one-time browser capture that feeds the graph for every later
caller. The two-call contract an agent sees is deliberately small: an \emph{intent
$\to$ ranked endpoint shortlist $\to$ execute} resolution returns the candidates
that answer an intent, and a single follow-up realizes the chosen endpoint call
against the live site.

\subsection{Capability layer: route discovery pipeline}
\paragraph{Traffic observation and filtering.} During a capture session the runtime
observes the requests a page makes and filters out noise (static assets, analytics,
telemetry), keeping the candidate first-party data endpoints.
\paragraph{Route extraction and normalisation.} Captured requests are normalised
into routes: method, URL template with typed path and query parameters, the auth
handshake, and a fingerprint of the response shape.
\paragraph{What the graph stores.} The graph stores callable routes and typed edges
(list$\to$detail, pagination, auth), plus a signed record of which routes were used
and what they returned --- value off-chain, commitment hash-chained.
\paragraph{Skill packaging.} Related routes for a domain are packaged into a
reusable skill that later agents install and execute without re-capturing.

\subsection{Composite scoring and intent resolution}
An intent is embedded and matched against the graph; candidates are ranked by a
composite of semantic similarity, observed reliability, freshness, and faithful
replay. The chosen route is executed; a content-addressed cache stores each
intermediate result under the hash of its content, so a second visit resolves
without launching a browser.

\subsection{Skill lifecycle and schema drift}
Routes decay --- schemas drift, auth flows mutate. The lifecycle tracks reliability
and demotes routes that stop reproducing, falling back to capture when a route goes
stale. The full accountability machinery that makes a freshness claim costly to
fake is the subject of the third paper~\cite{maintenance} and is not relitigated
here.

\section{Economic model}
\subsection{The adoption condition}
A cost-minimising agent will use the shared graph if and only if the fee to reuse a
mapped route undercuts the cost of rediscovering it from a browser:
\begin{equation}
  f_{\text{route}} \;<\; c_{\text{rediscovery}}. \label{eq:adopt}
\end{equation}
The rediscovery cost decomposes into the browser's latency, compute, and token
cost:
\begin{equation}
  c_{\text{rediscovery}} \;=\; c_{\text{latency}} + c_{\text{compute}} +
  c_{\text{tokens}}. \label{eq:redis}
\end{equation}
Because the mapping cost is paid once and amortised over every later caller,
\eqref{eq:adopt} holds for any route reused more than a handful of times.

\subsection{What agents pay for}
Discovery and internal-API routing are free; only execution against a priced route
settles. An \emph{unpaid \texttt{execute} gets HTTP 402}. Agents pay for the
value they consume --- a settled call against a maintained route --- not for the
act of looking.

\subsection{Why a shared commons can emerge}
Under \eqref{eq:adopt} the dominant strategy for each agent is to reuse rather than
re-map, and the dominant strategy for an indexer is to publish a route that others
will reuse and earn on. The commons emerges because cooperation out-competes
re-derivation for every party, not because participation is mandated.

\section{Route-level payment architecture}
All tiers are additive and collectively bounded by the adoption condition
\eqref{eq:adopt}.
\paragraph{Tier 1: skill installation.} A one-time fee to install a packaged skill,
priced well below a single browser rediscovery.
\paragraph{Tier 2: opt-in per-execution.} Per-call settlement on priced routes.
Paid execution \emph{settles over \textbf{x402}}, and the settlement \emph{splits
three ways} --- to the indexer who mapped the route, the platform, and (when
claimed) the site owner --- so \emph{x402 routes USDC to all three} in one
transaction. This \emph{three-way live split is shipped}.
\paragraph{Tier 3: search and routing fees.} An optional fee on a priced shortlist
returned by resolution.
\paragraph{Dynamic route pricing.} Prices track observed reliability and demand so
a maintained, reliable route earns more than a stale one, keeping the incentive
pointed at quality.

\section{Quality proofing and validation}
A reused route is only valuable if later agents can trust it. New submissions land
in a shadow state until corroborated by independent reuse; reliability is scored
from real outcomes (success ratio, consecutive failures, schema-drift signals); and
a route that stops reproducing is demoted. The cryptographic accountability that
makes a freshness claim \emph{bondable} is developed in the companion
papers~\cite{security,maintenance}.

\section{Implementation and evaluation}
We evaluate the API-native path against a browser-automation baseline (Playwright)
on a corpus of 94 live domains spanning search, listings, detail pages, and
authenticated reads. Table~\ref{tab:eval} reports the headline results. Both paths
make requests to the same target servers over the same network; the latency
difference is browser execution overhead.

\begin{table}[h]
\centering
\begin{tabular}{lr}
\toprule
Metric & Value \\
\midrule
Domains evaluated & 94 \\
Mean latency (Unbrowse) & 950\,ms \\
Mean latency (Playwright) & 3{,}404\,ms \\
Mean speedup & 3.6$\times$ \\
Median speedup & 5.4$\times$ \\
Best single-domain speedup & 30$\times$ \\
Mean cold-start (first capture) & 12{,}400\,ms \\
Tokens vs.\ browser baseline & 40$\times$ fewer \\
\bottomrule
\end{tabular}
\caption{API-native vs.\ browser-automation baseline across 94 live domains.}
\label{tab:eval}
\end{table}

The reuse mechanism is isolated in a runnable reference benchmark: a warm cached
replay categorically skips the cold mapping path, and the harness records a large
mean speedup on the reuse micro-benchmark (\texttt{reference/bench/bench\_reuse.py})
--- observed runs land in the $\approx$50--90$\times$ range, hardware-dependent, and
the gate asserts only that reuse is materially faster ($>2\times$) so the floor, not
a single headline number, is what is verified --- with a conservative
$\approx$27$\times$ observed gap on a single live URL
(\texttt{reference/bench/live-measured-1url.json}). The full
release-coverage methodology --- corpus shape, scoring rubric, and current numbers
--- is maintained alongside the product.

\paragraph{Cost savings.} The latency advantage is also a dollar advantage. A cold
browser rediscovery of a route costs an estimated \$0.10--0.53 in compute and tokens
(Table~\ref{tab:cost}); the shared graph turns that into a one-time
skill-installation fee of \$0.005--0.02, an order-of-magnitude reduction in the cost
of obtaining a working route, after which reuse is a content-addressed lookup. The
marketplace prices reuse below rediscovery by construction (Eq.~\ref{eq:adopt}), so
the saving is structural, not promotional.

\begin{table}[h]
\centering
\begin{tabular}{lr}
\toprule
Action & Estimated cost \\
\midrule
Browser rediscovery (cold, per route) & \$0.10 -- \$0.53 \\
Tier 1: skill installation (one-time) & \$0.005 -- \$0.02 \\
Reuse (content-addressed lookup) & $\approx$\$0 \\
\bottomrule
\end{tabular}
\caption{Per-route cost: browser rediscovery vs.\ shared-graph reuse.}
\label{tab:cost}
\end{table}

\paragraph{Coverage.} Speed and cost are worthless without coverage, so we measure
whether the route graph actually returns real results across a broad, diverse intent
space rather than on cherry-picked domains. A live coverage gate runs the real
\texttt{unbrowse search} binary over a diverse intent set and counts the fraction
that return real results; on the current graph it returns results on every probed
intent (coverage $=1.0$ against a $\geq 0.75$ release threshold). Coverage
is a gated, re-runnable measurement (\texttt{scripts/coverage-gate.sh}), not a static
claim --- it fails honestly when the graph cannot answer.

\subsection{The per-call figure is a floor, not the headline}
The 3.6$\times$ mean of Table~\ref{tab:eval} is a \emph{per-call} comparison: one
cached route execution against one cold browser rediscovery. A real agent task is
not one call. It is a sequence of browser interactions --- navigate, wait, click,
wait, fill, submit, paginate --- and each interaction costs latency, can fail and be
retried, and re-reads a growing page into the agent's context, so token cost
\emph{compounds} across the heavy DOM reads. The API-native path collapses the same
task into roughly one structured call: no multi-step navigation, no retry tax, no
compounding context. Accounting for this honestly, \emph{the per-call figure is a
lower bound, not the end-to-end advantage}: the per-task cost is the per-call cost
multiplied by the task's click count, its failure-retry factor, and its token
compounding. On a representative task (ten interactions, three DOM reads, an 80\%
per-step success rate) the model shows why multi-step browser tasks can exceed the
single-call speedup by a large margin. We ship this as a deterministic, runnable
reference (\texttt{reference/bench/task\_cost.py}) so the compounding is provable
arithmetic rather than a marketing multiplier.

\section{Discussion}
The first-party API is not always present (some surfaces are genuinely
render-bound), and a learned route can drift. Unbrowse treats both honestly: the
browser remains the bottom of a descent for the render-bound minority, and the
lifecycle demotes drifting routes. What the wedge buys is the common case in the
measured benchmark and in many modern web applications: a structured request the
site already makes, recovered once and shared. Two questions it deliberately leaves
open --- securing the descent below HTTP, and keeping the shared graph trustworthy
as it scales --- are the subjects of the two companion papers.

\section{Conclusion}
The graph is the asset. Every site an agent automates today is re-discovered from
scratch, per agent, per call --- the browser tax paid over and over. Turning that
one-off rediscovery into a shared, reusable route graph reduces that tax and makes
maintenance, not discovery, the place value compounds. Internal APIs carry a large
part of repeatable web-agent work when they are present and fresh. They are not
the whole stack: securing the layers below HTTP is \emph{Crypto Was All You
Needed}~\cite{security}, and keeping the graph fresh and bonded is the
\emph{Unbrowse Maintenance Network}~\cite{maintenance}. This paper is the wedge they
stand on.

\appendix
\section{Reproducibility}
The benchmark corpus, capture configuration, and scoring rubric are maintained with
the product; the reuse mechanism is reproduced in isolation by
\texttt{reference/bench/bench\_reuse.py} so the categorical
warm-vs-cold gap can be re-derived independently of any live endpoint.

\section{Benchmark domain list}
The 94 evaluated domains span search engines, marketplaces, social and news feeds,
developer platforms, and authenticated dashboards; the per-domain breakdown is
distributed with the release-coverage methodology.

\begin{thebibliography}{99}
\bibitem{security} L.~Tham. \emph{Crypto Was All You Needed: A Signed,
Layer-Descending Stack for Agent Computation}. Unbrowse AI, 2026.
\bibitem{maintenance} L.~Tham. \emph{Unbrowse Maintenance Network: Proof of
Indexing and Bonded Accountability in a Shared Route Graph}. Unbrowse AI, 2026.
\bibitem{webvoyager} H.~He et al. \emph{WebVoyager: Building an End-to-End Web Agent
with Large Multimodal Models}. arXiv:2401.13919, 2024.
\bibitem{mind2web} X.~Deng et al. \emph{Mind2Web: Towards a Generalist Agent for the
Web}. NeurIPS, 2023.
\bibitem{webarena} S.~Zhou et al. \emph{WebArena: A Realistic Web Environment for
Building Autonomous Agents}. ICLR, 2024.
\bibitem{toolformer} T.~Schick et al. \emph{Toolformer: Language Models Can Teach
Themselves to Use Tools}. NeurIPS, 2023.
\bibitem{gorilla} S.~G.~Patil et al. \emph{Gorilla: Large Language Model Connected
with Massive APIs}. arXiv:2305.15334, 2023.
\bibitem{mcp} Anthropic. \emph{Model Context Protocol}. 2024.
\url{https://modelcontextprotocol.io}.
\bibitem{x402} Coinbase. \emph{x402: An Open Protocol for Internet-Native Payments
over HTTP 402}. x402.org whitepaper, 2025. \url{https://github.com/coinbase/x402}.
\bibitem{ostrom} E.~Ostrom. \emph{Governing the Commons: The Evolution of
Institutions for Collective Action}. Cambridge University Press, 1990.
\end{thebibliography}

\end{document}
